\section{Government Agent: de la Task Demo al Workflow System}

La última pregunta de toda la clase es la más cercana a la agent engineering actual. Alswaha distingue entre la AI enhancement del legacy B2B SaaS y el AI-native player, y sostiene que la dificultad no está solo en el modelo, sino en los data, el business model y el human-in-the-loop. Dicho con más precisión, el task-level success solo se convierte en un sistema de servicio público si cruza tres puertas: workflow integration, multi-agent coordination y public accountability.

\subsection{Task, workflow y multi-agent: tres expansiones del estado}

Un \term{task-level agent} completa drafting, classification o lookup sobre entradas locales; un \term{workflow-level agent} hace avanzar un case a través de varios sistemas, permisos y estados de larga duración; un \term{multi-agent system} hace que varios agents se repartan el trabajo, compartan memory, coordinen herramientas y resuelvan conflictos. Cada expansión aumenta el state space, la failure interaction y la audit complexity.

\lecturefigure{16-government-agent-maturity.png}{Un Agent gubernamental debe cruzar, en orden, task, workflow, multi-agent y public accountability.}{Subtítulos locales 49:50--52:17; redibujo conceptual.}

\begin{knowledgebox}{Lectura de la figura: por qué la task demo es la que con más facilidad tiene éxito}
En una task demo, las entradas pueden filtrarse manualmente, el estado es corto, los permisos son pocos y las fallas pueden ignorarse; un workflow real se topa con campos faltantes, cases duplicados, ownership entre departamentos, deadlines, appeals y policy exceptions. Multi-agent introduce además shared memory, conflicting plans, deadlock y cascading tool error; por ello, la tasa de éxito en producción no puede extrapolarse a partir de un benchmark de una sola ronda.
\end{knowledgebox}

\subsection{Data layer, human-in-the-loop y business model}

\term{human-in-the-loop} (humano en el ciclo) no significa «que una persona confirme con un clic cualquier resultado», sino turnar los casos de alto riesgo, los de baja confianza, las excepciones de política y las vías de apelación a una persona con la autoridad correspondiente, y recopilar sus razones como evidencia para mejorar. La data layer requiere canonical entity, lineage, quality rule, permission y retention; el business model, por su parte, responde quién paga el despliegue, los errores, la revisión humana y la actualización continua.

\begin{importantbox}{Matriz de aceptación de un Agent de servicio público}
\begin{tabularx}{\textwidth}{L{0.18\textwidth}L{0.29\textwidth}X}
\toprule
Dimensión & Pregunta clave & Evidencia medible \\
\midrule
Calidad de la tarea & ¿La salida es correcta, completa y explicable? & expert review, error taxonomy, calibration \\
Integridad del proceso & ¿Hace avanzar correctamente el state y las exception? & end-to-end case success, rework, deadline miss \\
Equidad y derechos & ¿Se puede apelar? ¿Produce un impacto diferenciado entre grupos? & subgroup audit, appeal outcome, override reason \\
Seguridad y privacidad & ¿Los permisos, los datos y las tool están minimizados? & access log, red-team, leakage/abuse test \\
Viabilidad económica & ¿La automatización reduce de verdad el total cost / cycle time? & compute + integration + human review + incident cost \\
\bottomrule
\end{tabularx}
\end{importantbox}

\begin{warningbox}{«El modelo no es la dificultad» tampoco debe tomarse como absoluto}
Los data y el workflow suelen ser el principal cuello de botella, pero la language coverage, el reasoning, la tool reliability, la calibration y la robustness del base model siguen limitando el sistema. La conclusión más precisa es esta: la capacidad del modelo es una condición necesaria, y el éxito en producción lo determina el producto de model, data, process, people, policy y economics.
\end{warningbox}

\teachervoice{La teacher voice más transferible de toda la clase aparece al final: algunos equipos conectan los vertical data directamente al modelo y solo obtienen task augmentation; un sistema verdaderamente AI-native rediseña la data creation, el human feedback y la workflow ownership.}

\subsection{Resumen de la sección}

El límite de un Government Agent no está en la interfaz de chat, sino en el case state, los permisos, los datos, las excepciones, las apelaciones y la responsabilidad. Una task demo puede estar impulsada por el modelo; un workflow system debe estar impulsado por una arquitectura sociotécnica completa.
